\chapter{グローバル・マクロとキャリー}
\label{ch:macro}

\section{裁量のマクロとシステマティックなマクロ}

グローバル・マクロは、経済の見通しから、金利・通貨・株価指数・商品の方向に賭ける。有名な運用者の多くは裁量であり、その判断の過程を論文で検証することは難しい。研究で扱えるのは、規則で書ける部分、つまり\textbf{キャリー}、\textbf{割安}（購買力平価などからの乖離）、\textbf{トレンド}（第 \ref{ch:trend} 章）である。この章は、主にキャリーを扱う。

\section{通貨のキャリー}

\subsection{金利平価の破れ}

金利の高い通貨の預金は、金利の低い通貨の預金より多くの利息を生む。カバーなしの金利平価（UIP）が成り立つなら、その差は、高金利の通貨の将来の値下がりでちょうど打ち消され、平均的な利益はゼロになるはずである。

\begin{paperbox}[Fama（1984）\cite{fama1984}]
先渡しの為替レートと、その後の直物の為替レートの関係を検証し、UIP が成り立たないことを示した。金利の高い通貨は、UIP の予測ほど値下がりせず、平均ではむしろ値上がりする傾向がある（フォワード・プレミアムのパズル）。
\end{paperbox}

そこで、金利の低い通貨（円、スイス・フランなど）を借りて売り、金利の高い通貨を買う\textbf{キャリー・トレード}は、平均すると利益を上げてきた。

\subsection{なぜ儲かるのか：リスクの対価}

\begin{paperbox}[Lustig・Roussanov・Verdelhan（2011）\cite{lrv2011}]
通貨を金利の水準で並べたポートフォリオのリターンが、二つの共通のファクターでよく説明できることを示した。一つはドルの全体の水準、もう一つは高金利の通貨を買い低金利の通貨を売るキャリーのファクター（HML\textsubscript{FX}）である。高金利の通貨は、世界の株式市場の変動率が上がる局面（世界的な危機）で損をする。キャリーの収益は、この世界共通のリスクを引き受ける対価だと解釈した。
\end{paperbox}

Menkhoff・Sarno・Schmeling・Schrimpf（2012）\cite{msss2012} は、世界の為替の変動率の予想外の上昇が、キャリーのリターンの横断面を説明することを示した。高金利の通貨は、為替の変動率が急に上がるときに損をする。

\subsection{急落}

\begin{paperbox}[Brunnermeier・Nagel・Pedersen（2008）「キャリー・トレードと通貨の急落」\cite{bnp2008}]
キャリー・トレードのリターンは負の歪度を持つ（ふだんは少しずつ儲け、まれに急に大きく損をする）。金利差が大きいほど、その後の急落のリスクも大きい。市場の変動率（VIX）が上がると、投機筋がキャリーのポジションを解消し、高金利の通貨が急落する。資金の流動性の制約が、急落の引き金になる（\chapref{ch:risk} の流動性のらせん）。
\end{paperbox}

2008 年の秋には、円の急騰で多くのキャリー・トレードが数週間のうちに大きく損をした。「蒸気ローラーの前で小銭を拾う」とたとえられる構造である。

\section{あらゆる資産のキャリー}

\begin{paperbox}[Koijen・Moskowitz・Pedersen・Vrugt（2018）「キャリー」\cite{kmpv2018}]
キャリーを「価格が変わらなかったときに得られる収益」と一般化し、株価指数（配当利回りと先物の価格の関係）、国債（利回り曲線の傾き）、商品（期近と期先の価格差）、通貨（金利差）、オプションなどで定義した。どの資産クラスでも、キャリーの高い資産が低い資産より高いリターンを持つ。資産クラスをまたいだキャリーの戦略の間の相関は低く、分散した「世界のキャリーのファクター」は高いリスク調整後のリターンを持つ。一方で、世界の景気後退や流動性の低下の局面では損をする。
\end{paperbox}

キャリーは、市場の「予想」（先物の価格）と、「何も変わらない」という予測のずれに賭ける戦略と見ることもできる。

\section{危機の局面}

Ang・Bekaert（2002）\cite{angbekaert2002} は、国際的な株式の資産配分を、「平時」と「相関と変動率が高い局面」を確率的に行き来するモデル（レジーム・スイッチング、\chapref{ch:model}）で分析した。下げ相場の局面では国の間の相関が上がり、国際分散の効果が小さくなるが、それでも分散の価値は残ることを示した。マクロの運用者が、局面の変化を見越して資産配分を動かす考え方の、理論的な裏付けの一つである。

\section{裁量のマクロの成績}

裁量のマクロのファンドの成績は、ヘッジファンドのデータベース（第 \ref{ch:hfperf} 章）で調べるしかない。Deep Research の答え 1 は、マクロのファンドのシャープレシオが上げ相場・下げ相場・危機でそれぞれ 1.0・0.8・0.9 だったという数値を挙げたが、出典を確かめられなかったので使わない。データベースの偏り（第 \ref{ch:hfperf} 章）を考えると、裁量のマクロの「腕」の大きさについて、確かなことは言いにくい。

\section{個人や研究で再現できるか}

\begin{itemize}
\item 各国の短期金利・為替レート・先物の価格は、公開のデータで入手しやすい。通貨のキャリーや、資産クラスをまたぐキャリーの戦略の検証は、比較的容易である。
\item ただし、キャリーの収益は、まれな急落の損失で大きく削られる。検証の期間に大きな急落が含まれていないと、成績が過大に見える。
\item 実運用では、個人の為替の取引のスプレッドや、借入れの費用が、機関投資家より不利である。
\end{itemize}
